%%%%%%%%%%  scan idx 602  |  folio 589  %%%%%%%%%%
\SGAPrintPage{589}%
and the composite of these homomorphisms for \(0\leq i\leq n-1\) defines a homomorphism (7.17.4.1), which one can show (by proceeding as in 7.3) to be the same as the one defined above by means of a blow-up. The advantage of the second method, however, is that here one even has homomorphisms of the form (7.16.1):
\[
  K_\bullet(X_{\bar t_{i+1}})\longrightarrow K_\bullet(X_{\bar t_i}),
\]
(with suitable coherent choices of the algebraic closures \(k(\bar t_i)\) of the \(k(t_i)\)), which allow one to define by composition a homomorphism
\sgaStableTarget{sga6.e10.eq.7-17-4-2}\begin{equation}\tag{7.17.4.2}
  K_\bullet(X_{\bar t})\longrightarrow K_\bullet(X_{\bar s}),
\end{equation}
that probably depends essentially on the choice of the \(S_i\), that is, of the \(t_i\).

\sgaStableTarget{sga6.e10.para.7-17-5}\textbf{7.17.5.} The construction of a specialization homomorphism (7.17.4.2) can be generalized to the case in which one no longer assumes that \(S\) is regular at \(s\) and that \(t\) is a maximal point, provided that resolution of singularities is available for the integral closed subscheme \(T\) of \(S\) with generic point \(t\), or merely for its localization at \(s\); this hypothesis holds in particular when \(\underline{O}_{S,s}\) is excellent of characteristic zero, by Hironaka. Indeed, one can then find a morphism \(S'\to \Spec(\underline{O}_{T,s})\) with \(S'\) regular, and, taking a closed point of the fiber \(S'_s\), apply the preceding construction to \(X'=X\times_S S'\) over \(S'\), the generic point \(t'\) of \(S'\) and its closed point \(s'\). This yields a homomorphism of the form (7.17.4.2), since \(k(s')\) is a finite extension of \(k(s)\) and therefore has the same algebraic closure.

\sgaStableTarget{sga6.e10.para.7-17-6}\textbf{7.17.6.} In all the constructions considered in section 7.17, the construction of a specialization homomorphism has been reduced to the

%%%%%%%%%%  scan idx 603  |  folio 590  %%%%%%%%%%
\SGAPrintPage{590}%
case where the base is a trait. It follows that the properties of the specialization homomorphism reviewed in the preceding sections, in the case where the base is a trait, extend automatically to the case of a more general base. Of course, the methods of definition considered apply equally well to \(\Gr_\bullet\), \(K^\bullet\) (when \(f\) is smooth), and \(\Gr^\bullet\) (when, in addition, \(f\) is separated).

\sgaStableTarget{sga6.e10.para.7-17-7}\textbf{7.17.7.} Finally, note that the definition of a homomorphism (7.17.3.1) or (7.17.3.2), or of a homomorphism (7.17.4.2), depends essentially on certain arbitrary choices (a suitable trait over \(S\) in 7.17.3, a sequence of \((S_i)\) in 7.17.4, and, in addition, a resolution of singularities in 7.17.5), and likewise for the variants for \(\Gr_\bullet\), etc. Obvious examples can be constructed with \(s\), say, an ordinary double point of an algebraic curve and \(X\) a projective smooth scheme over \(S\) (for example, an abelian scheme of relative dimension 2), in which the specialization homomorphism on \(\Gr^1=\Pic\) and even on \(\Gr^1_{\mathrm{alg}}=\mathrm{NS}\) (the Néron--Sévéri group) actually depends on the choice of the trait \(S'\) chosen over \(S\) or, what amounts to the same thing here, on the choice of a point over \(s\) in the normalization of \(S\).

\sgaStableTarget{sga6.e10.sec.7.18}\subsection*{7.18. Appendix to the Appendix: extension of the base field for certain functors with values in commutative groups}
We spell out here a technical point raised in 7.17.4.

\sgaStableTarget{sga6.e10.para.7-18-1}\textbf{7.18.1.} Let \(k\) be a field, \(C\) the category of integral \(k\)-algebras and flat homomorphisms of such algebras, and \(F\) a functor on \(C\) with values in the category \((\mathrm{Ens})\) of sets. Suppose that \(F\) commutes with filtered inductive limits. Let \(K\) be an extension of \(k\); we seek

%%%%%%%%%%  scan idx 604  |  folio 591  %%%%%%%%%%
\SGAPrintPage{591}%
criteria for \(F(k)\longrightarrow F(K)\) to be injective (resp. surjective).

\sgaStableTarget{sga6.e10.para.7-18-2}\textbf{7.18.2.} Suppose that we are in one of the following two cases:
\begin{enumerate}[label=\alph*)]
\item The field \(k\) is separably closed, and \(K/k\) is a separable extension.
\item The field \(k\) is infinite, and \(K\) is a pure extension of it.
\end{enumerate}
Then the map \(F(k)\to F(K)\) is injective.

Indeed, using the fact that \(K\) is the filtered inductive limit of subalgebras of finite type over \(k\), with flat transition morphisms, and the commutation of \(F\) with filtered inductive limits, one is reduced in case a) to proving that, if \(A\) is an integral algebra of finite type over \(k\) whose fraction field \(K\) is separable over \(k\), then \(F(k)\to F(A)\) is injective. Now the hypotheses on \(k\) and on \(K\) imply that there exists a \(k\)-homomorphism \(A\to k\); the composite \(F(k)\to F(A)\to F(k)\) is then the identity, whence the conclusion. In case b), passage to the limit similarly reduces the question to proving that, if \(A\) is an integral algebra of finite type over \(k\) whose fraction field is a pure extension of \(k\), then \(F(k)\to F(A)\) is injective. Since \(k\) is infinite, the hypothesis again implies that there exists a \(k\)-homomorphism \(A\to k\), and one concludes as above.

\sgaStableTarget{sga6.e10.para.7-18-3}\textbf{7.18.3.} Now suppose that \(F\) is a functor with values in the category \((\mathrm{Ab})\) of abelian groups, and suppose that for every finite extension \(k'\) of \(k\) a norm homomorphism
\[
  N_{k'/k}:F(k')\longrightarrow F(k)
\]
has been defined such that the composite
\[
  F(k)\longrightarrow F(k')\xrightarrow{N_{k'/k}}F(k)
\]
is multiplication by \([k':k]\). Under these conditions:
\begin{enumerate}[label=(\roman*)]

%%%%%%%%%%  scan idx 605  |  folio 592  %%%%%%%%%%
\item
\SGAPrintPage{592}%
For every extension \(K\) of \(k\), the kernel of \(F(k)\to F(K)\) is a torsion group.
\item For every pure extension \(K\) of \(k\), the homomorphism \(F(k)\to F(K)\) is injective. (NB. Note that \(k\) is not necessarily assumed infinite).
\end{enumerate}

In case (i), one is reduced to proving that every element \(x\) in the kernel of \(F(k)\to F(A)\), where \(A\) is an integral algebra of finite type over \(k\), is a torsion element. Using a \(k\)-homomorphism \(A\to k'\), where \(k'\) is a finite \(k\)-algebra, one is reduced to the case where \(A=k'\) is a finite extension of \(k\). One then has \([k':k]x=0\) by the property postulated above for the norm homomorphism. This proves that \(x\) is a torsion element. One even has \(x=0\) when the greatest common divisor of the degrees of the finite extensions \(k'\) of \(k\) that are quotients of \(A\) equals 1. This is the case when the fraction field \(K\) of \(A\) is a pure extension of \(k\), as is readily checked (only the case where \(k\) is finite requires a new verification). This observation then gives the conclusion of (ii).

\sgaStableTarget{sga6.e10.para.7-18-4}\textbf{7.18.4.} Suppose that, for every nonempty affine open subset \(U\) of the affine space \(\Spec k[T_1,\ldots,T_n]\), with \(A\) the ring of \(U\), the homomorphism \(F(k)\to F(A)\) is surjective (compare IX 1.6). Then, using the commutation of \(F\) with filtered inductive limits, one finds that for every pure extension \(K\) of \(k\), \(F(k)\to F(K)\) is surjective. If one also assumes that \(k\) is infinite, or that norm homomorphisms can be defined as in 7.18.3, one concludes that this homomorphism is even bijective.

\sgaStableTarget{sga6.e10.para.7-18-5}\textbf{7.18.5.} As an example, consider a scheme \(X\) of finite type over \(k\), and set
\[
  F(A)=K_\bullet(X\otimes_k A),
\]

%%%%%%%%%%  scan idx 606  |  folio 593  %%%%%%%%%%
\SGAPrintPage{593}%
which is a functor on the category \(C\) (by (IV 2.12)). By (IV 3.2.3), this functor commutes with filtered inductive limits. Moreover, when \(k'\) is a finite extension of \(k\), the projection \(p:X\otimes_k k'\to X\) is a finite, flat morphism of finite presentation, and therefore defines a direct-image homomorphism \(p=N_{k'/k}:K_\bullet(X\otimes_k k')\longrightarrow K_\bullet(X)\). One immediately verifies that the condition on the norm homomorphisms considered in 7.18.3 does hold. Finally, if \(A\) is the coordinate algebra of an open subset \(U\) of the affine space \(\Spec k[T_1,\ldots,T_n]\), then \(F(k)\to F(A)\) is identified with the map \(K_\bullet(X)\longrightarrow K_\bullet(X\times_k U)\), which is surjective for \(X\) noetherian, by IX 1.6. Using the results proved above, one then obtains:
\begin{enumerate}[label=(\roman*)]
\item If \(K/k\) is a pure extension of \(k\), then \(K_\bullet(X)\longrightarrow K_\bullet(X_K)\) is injective, and is even bijective when \(X\) is noetherian.
\item If \(k\) is separably closed and \(K/k\) is a separable extension, then \(K_\bullet(X)\longrightarrow K_\bullet(X_K)\) is injective.
\item For every extension \(K\) of \(k\), the kernel of \(K_\bullet(X)\longrightarrow K_\bullet(X_K)\) is a torsion group.
\end{enumerate}
These results remain valid when one assumes only that \(X\) is noetherian, provided one restricts to extensions \(K\) of \(k\) that are of finite type (so that \(X_K\) is still noetherian). Moreover, essentially the same results hold for \(K^\bullet\) in place of the functor \(K_\bullet\), with the sole difference that in (i) one cannot assert bijectivity in general, for lack of the homotopy theorem for \(K^\bullet\) in place of \(K_\bullet\); to conclude bijectivity, one would have to assume that \(X\) is noetherian and regular, which implies that the groups \(K^\bullet\) and \(K_\bullet\) are the same for \(X\) and for \(X_K\), respectively.

